\documentclass{ltjsarticle}
\usepackage{luatexja}
\usepackage[haranoaji]{luatexja-preset}

% Add because main language is not English
\usepackage{enumitem}
\setlist[itemize]{label=-}

% --- PACKAGES FOR MATH & DIAGRAMS ---
\usepackage{amsmath}
\usepackage{amssymb}
\usepackage{amsthm}
\usepackage{mathtools}
\usepackage{tikz}
\usepackage{tikz-cd}
\usetikzlibrary{shapes, positioning, arrows.meta, decorations.pathreplacing, backgrounds}

% --- THEOREM ENVIRONMENTS ---
\theoremstyle{definition}
\newtheorem{definition}{定義}[section]
\newtheorem{example}[definition]{例}
\theoremstyle{plain}
\newtheorem{theorem}[definition]{定理}
\newtheorem{lemma}[definition]{補題}
\newtheorem{proposition}[definition]{命題}
\newtheorem{remark}[definition]{注釈}

% --- METADATA ---
\title{\textbf{Cat\_Dの位相幾何学的応用：構造塔による空間の階層化}\\
\large Lean 4による実装 \texttt{P3\_Topological.lean} の数学的解説}
\author{
    \textbf{TeX生成}: Gemini (Google) \\
    \textbf{Lean実装}: Codex (OpenAI)
}
\date{\today}

\begin{document}

\maketitle

\begin{abstract}
本稿では、形式化数学ライブラリ \texttt{MyProjects.ST.CAT.Cat\_D} における位相空間論への応用について解説する。特に、第二可算空間における開集合の複雑性に基づく階層構造（Open Set Tower）と、閉包演算の反復に基づく階層構造（Closure Tower）に焦点を当て、連続写像がこれらの階層構造の間にもたらす射の性質について、提供されたLeanコードに基づき数学的に定式化する。
\end{abstract}

\tableofcontents

\section{はじめに}
提供されたLeanコード \texttt{P3\_Topological.lean} は、抽象的な階層構造（Tower）を扱うカテゴリ \texttt{Cat\_D} を、具体的な位相空間論に応用したものである。
主なアイデアは、位相空間 $X$ の部分集合族（開集合や閉集合など）に対し、ある種の「複雑さ」や「生成ステップ」を表す自然数 $n \in \mathbb{N}$ を割り当て、その構造を保存する写像を研究することにある。

\section{開集合の階層 (Open Set Tower)}

\subsection{定義と直観}
第二可算空間（Second Countable Space）、あるいはより一般に可算基底を持つ位相空間において、任意の開集合は基底の和集合として表される。計算機科学的な観点、あるいは有限近似の観点からは、「いくつの基底を使えばその開集合を表現できるか」は重要な指標となる。

\begin{definition}[有限基底和]
位相空間 $X$ とその開基底 $\mathcal{B}$ を考える。開集合 $U \subseteq X$ が \textbf{レベル $n$ の有限基底和} であるとは、以下を満たすことである。
\[
\exists \mathcal{S} \subseteq \mathcal{B}, \quad |\mathcal{S}| \le n \quad \text{かつ} \quad U = \bigcup_{V \in \mathcal{S}} V
\]
Leanコード内では \texttt{IsFiniteUnionOfBasis} として定義されている。
\end{definition}

\begin{definition}[開集合構造塔]
$(X, \mathcal{B})$ に対する \textbf{開集合構造塔 (Open Set Tower)} とは、以下の階層構造を持つオブジェクトである。
\begin{itemize}
    \item \textbf{台集合}: $\text{Carrier} = \{ U \subseteq X \mid U \text{ is Open} \}$
    \item \textbf{層 (Layer)}: 第 $n$ 層 $L_n$ は、高々 $n$ 個の基底の和で表される開集合全体。
\end{itemize}
\[
L_n = \{ U \in \text{Open}(X) \mid \exists \{V_1, \dots, V_k\} \subseteq \mathcal{B}, k \le n, U = \bigcup_{i=1}^k V_i \}
\]
\end{definition}

この定義により、空集合は $L_0$ に（0個の和）、基本開集合自体は $L_1$ に属することになる。

\begin{figure}[htbp]
\centering
\begin{tikzpicture}[scale=0.8]
    % Draw the tower
    \draw[thick, fill=blue!5] (-4,0) -- (4,0) -- (3,6) -- (-3,6) -- cycle;
    
    % Layers
    \draw[dashed] (-3.8, 1) -- (3.8, 1);
    \draw[dashed] (-3.6, 2) -- (3.6, 2);
    \draw[dashed] (-3.4, 3) -- (3.4, 3);
    
    % Labels
    \node at (0, 0.5) {$L_0 = \{\emptyset\}$};
    \node at (0, 1.5) {$L_1 = \mathcal{B}$ (基底)};
    \node at (0, 2.5) {$L_2 = \{V_1 \cup V_2 \mid V_i \in \mathcal{B}\}$};
    \node at (0, 4.5) {$\vdots$};
    \node at (0, 5.5) {$L_n$ (複雑さ $n$)};
    
    \node[right] at (4, 3) {\textbf{Open Set Tower}};
\end{tikzpicture}
\caption{開集合の階層構造のイメージ}
\end{figure}

\subsection{連続写像と一様有界性}

位相空間の間の連続写像 $f: X \to Y$ は開集合の逆像を開集合に写すが、階層構造（基底の個数）を保つとは限らない。Leanコードでは、この問題に対して「一様上界 (Uniform Bound)」という概念を導入して解決している。

\begin{definition}[逆像の基底被覆数に関する一様上界]
$X, Y$ をそれぞれ基底 $\mathcal{B}_X, \mathcal{B}_Y$ を持つ空間とし、$f: X \to Y$ を連続写像とする。
ある定数 $k \in \mathbb{N}$ が存在して、任意の $V \in \mathcal{B}_Y$ に対し、$f^{-1}(V)$ が $\mathcal{B}_X$ の高々 $k$ 個の元の和で表されるとき、$f$ は \textbf{基底に関して $k$-有界} であるという。
Leanコード内の \texttt{PreimageBasisBound} に対応する。
\end{definition}

\begin{theorem}[積の法則]
もし $f$ が基底に関して $k$-有界であるならば、任意の $U \in L_n(Y)$ に対して
\[
f^{-1}(U) \in L_{n \times k}(X)
\]
が成り立つ。すなわち、複雑さは高々 $k$ 倍にしかならない。
\end{theorem}

\begin{proof}[証明の概略]
$U \in L_n(Y)$ とする。定義より $U = \bigcup_{i=1}^m V_i$ ($V_i \in \mathcal{B}_Y, m \le n$) と書ける。
逆像をとると、
\[
f^{-1}(U) = f^{-1}\left(\bigcup_{i=1}^m V_i\right) = \bigcup_{i=1}^m f^{-1}(V_i)
\]
仮定より、各 $f^{-1}(V_i)$ は高々 $k$ 個の $\mathcal{B}_X$ の元の和である。
したがって、全体として $f^{-1}(U)$ は高々 $m \times k \le n \times k$ 個の基底の和となる。
\end{proof}

この定理により、連続写像 $f$ は \texttt{Cat\_D} における射（Tower Morphism）を誘導する。

\begin{figure}[htbp]
\centering
\begin{tikzpicture}
    % Space Y
    \node (Y) at (0,0) [cylinder, shape border rotate=90, draw, minimum height=4cm, minimum width=2cm, fill=yellow!10] {};
    \node at (0,-2.5) {Space $Y$ (Tower $T_Y$)};
    \node at (0, 1) {$U \in L_n$};
    
    % Space X
    \node (X) at (8,0) [cylinder, shape border rotate=90, draw, minimum height=4cm, minimum width=2cm, fill=blue!10] {};
    \node at (8,-2.5) {Space $X$ (Tower $T_X$)};
    \node at (8, 1) {$f^{-1}(U) \in L_{n \times k}$};
    
    % Map
    \draw[->, thick] (Y) -- (X) node[midway, above] {$f^{-1}$ (Map)};
    \draw[->, thick, dashed, bend right] (0.5, 1) to node[below] {複雑さの拡大 $\times k$} (7.5, 1);

    % Basis expansion visual
    \node[draw, rectangle, fill=white, inner sep=2pt] at (0, -1) {\small $V \in \mathcal{B}_Y$};
    \node[draw, rectangle, fill=white, inner sep=2pt] at (8, -1) {\small $\bigcup_{j=1}^k W_j, W_j \in \mathcal{B}_X$};
    
    % Mapping arrow for basis
    \draw[|->] (1.2, -1) -- (5.5, -1);
\end{tikzpicture}
\caption{写像による階層の移動と係数 $k$}
\end{figure}

\section{閉包演算の階層 (Closure Tower)}

\subsection{閉包の反復}
もう一つの応用例として、閉包演算（Closure）の反復による階層化が実装されている。
通常、位相空間の部分集合 $A$ に対し、$\text{cl}(A) = A$ （閉集合）であれば演算は停止するが、一般の集合に対しては閉包をとる操作によって集合が変化する。

Leanコードでは、以下の反復列を考えている。
\begin{align*}
    A^{(0)} &= A \\
    A^{(i+1)} &= \text{cl}(A^{(i)})
\end{align*}

\begin{definition}[閉包反復レベル]
集合 $A$ の閉包反復レベル (Closure Iteration Level) とは、この反復列が安定する（あるいは閉集合になる）までの最小ステップ数である。
\end{definition}

実装されている \texttt{closureTower} では、教育的な簡易版として以下の2層構造が示唆されている（完全版では順序数やCantor-Bendixsonランクへの拡張が視野にある）。
\begin{itemize}
    \item $L_0$: 既に閉集合であるもの（$\text{cl}(A) = A$）。
    \item $L_n$: $n$ 回の閉包操作で安定するもの。
\end{itemize}

連続写像 $f: X \to Y$ は $f(\text{cl}(A)) \subseteq \text{cl}(f(A))$ という性質を持つため、閉包構造塔の間の射についても考察されている。

\section{結論}
\texttt{P3\_Topological.lean} は、位相空間論における「基底の個数」や「閉包の回数」といった離散的な数値を \texttt{Cat\_D} の構造塔として捉え直す試みである。
特に \texttt{IsFiniteUnionOfBasis.preimage\_mul} （補題：有限基底和の逆像）は、連続写像がこの構造塔の射となるための十分条件（一様有界性）を明確にしており、計算機科学的なリソース制約付き位相幾何学の基礎をなすものである。

\end{document}